% ---- pg-592 （书页 580） ----

\noindent 12.2\quad 用\hyperref[tab:12.1]{表 12.1} 中的相似性系数 5,6 和 7 重复练习 12.1.

\noindent 12.3\quad 证明: 对具有后面所列频数的二值变量, 样本相关系数[见\hyperref[eq:12-11]{式(12-11)}]可以写为
\[
r=\frac{ad-bc}{[(a+b)(a+c)(b+d)(c+d)]^{1/2}}
\]
二值变量的频数为

\begin{mvtable}
\setlength{\tabcolsep}{10pt}%
\begin{tabular}{cc|cc}
\toprule
& & \multicolumn{2}{c}{变量 2} \\
\cmidrule(lr){3-4}
& & 0 & 1 \\
\midrule
\multirow{2}{*}{变量 1} & 0 & $a$ & $b$ \\
& 1 & $c$ & $d$ \\
\bottomrule
\end{tabular}
\end{mvtable}

\noindent 12.4\quad 试证\hyperref[tab:12.1]{表 12.1} 中的相似性系数 1,2,3 之间有单调关系.

提示: $(b+c)=p-(a+d)$. 故有(比如说)
\[
\frac{a+d}{a+d+2(b+c)}=\frac{1}{1+2[p/(a+d)-1]}
\]
这个方程表达了系数 3 与 1 之间的关系. 求其他各对系数的类似表示.

\noindent 12.5\quad 考虑距离矩阵
\[
\begin{array}{r@{\;}c}
 & \begin{array}{cccc}1&2&3&4\end{array}\\[2pt]
\begin{array}{r}1\\2\\3\\4\end{array} &
\left[\begin{array}{cccc}
0&&&\\
1&0&&\\
11&2&0&\\
5&3&4&0
\end{array}\right]
\end{array}
\]
用下述方法对以上 4 个项目进行聚类.

\noindent (a) 单连接分层方法.

\noindent (b) 完全连接分层方法.

\noindent (c) 平均连接分层方法.

画出树图, 并比较以上三种方法所得到的结果.

\noindent 12.6\quad 五个项目间的距离阵如下:
\[
\begin{array}{r@{\;}c}
 & \begin{array}{ccccc}1&2&3&4&5\end{array}\\[2pt]
\begin{array}{r}1\\2\\3\\4\\5\end{array} &
\left[\begin{array}{ccccc}
0&&&&\\
4&0&&&\\
6&9&0&&\\
1&7&10&0&\\
6&3&5&8&0
\end{array}\right]
\end{array}
\]
用练习 12.5 中的三种方法对这五个项目进行聚类. 画出树图并比较结果.

\noindent 12.7\quad 例 8.5 中曾给出五种股票间的相关系数. 这些相关系数在保留两位小数后重新列于下表:
